\section*{ABSTRACT}
\addcontentsline{toc}{section}{ABSTRACT}

A vehicle access management system was developed for a parking facility in Oruro,
Bolivia, to replace the manual registration of entries and exits with an automated
process based on computer vision and deep learning. The system relied on the
facility's existing cameras and integrated a dedicated dataset, detection and
classification models, a license plate recognizer, and a web platform with a
database, a task queue, and real-time notification. A total of 640 real license
plates were verified and 360,000 synthetic samples were generated. Vehicle
detection reached 0.901 mAP@50; plate detection, 0.9714 recall; vehicle type
classification, 0.8887 F1; and color classification, 0.8124 macro F1. License
plate recognition obtained 0.851 character-level accuracy and 52.6\% per-crop
exact match, both limited by capture resolution. Two criteria fell below the
target threshold (color F1 and character accuracy), event determination
remained inconclusive because of sample size, and per-event latency was not
measured; the research hypothesis is therefore rejected and the null hypothesis is
accepted under the predefined criterion. The cost estimate yielded 62 function
points and a total cost of 33,170 bolivianos. The system is built and integrated
by components, but it does not meet the set of performance criteria in its current
state.

\vspace{1em}
\noindent\textbf{Keywords:} computer vision, deep learning, license plate
recognition, access management, parking.
